\subsection{The preparation theorem}

In this no. A will denote a local ring, m its maximal ideal and k = A/m its residue field. Suppose that A is Hausdorff and complete with the m-adic topology. Let B = A{[}{[}X{]}{]}; it is local ring whose maximal ideal N is generated by m and X; with the \(\mathfrak{N}\) -adic topology, B is Hausdorff and complete (Chapter III, § 2, no. 6, Proposition 6).

For every formal power series

\[
f = \sum_ {i = 0} ^ {\infty} a _ {i} \mathbf {X} ^ {i} \in \mathbf {B},
\]

we write

\[
\bar {f} = \sum_ {i = 0} ^ {\infty} \bar {a} _ {i} X ^ {i} \in k [ [ X ] ],
\]

where \(\bar{a}_{i}\) denotes the canonical image of \(a_{i}\) in k. The series \(\bar{f}\) will be called the reduced series of f; if \(\bar{f} \neq 0\) , the order of \(\bar{f}\) (that is the least integer s such that \(a_{s} \notin m\) ) will be called the reduced order of f.

PROPOSITION 5. Let \(f \in \mathbf{B}\) be a series whose reduced series is non-zero. Let \(s\) denote its reduced order and \(\mathbf{M}\) the sub-A-module of \(\mathbf{B}\) with basis \(\{1, X, \ldots, X^{s-1}\}\) . Then \(\mathbf{B}\) is the direct sum of \(\mathbf{M}\) and \(f\mathbf{B}\) and \(f\) is not a divisor of zero in \(\mathbf{B}\) .

\begin{enumerate}
\def\labelenumi{(\alph{enumi})}
\tightlist
\item
  We show that \(f\mathbf{B} \cap \mathbf{M} = (0)\) . Suppose that there is a relation:
\end{enumerate}

\[
\left(\sum_ {i = 0} ^ {\infty} b _ {i} \mathrm{X} ^ {i}\right) f = r _ {0} + r _ {1} \mathrm{X} + \dots + r _ {s - 1} \mathrm{X} ^ {s - 1} \quad (b _ {i} \in \mathrm{A}, r _ {j} \in \mathrm{A}).\tag{3}
\]

We show that the \(b_{i}\) (and hence the \(r_{j}\) ) are all zero, which will prove in particular that \(f\) is not a divisor of zero in B. Since A is Hausdorff, it suffices to show that \(b_{i} \in \mathfrak{m}^{n}\) for all \(i \geqslant 0\) and all \(n \geqslant 0\) . It is obvious for \(n = 0\) . We shall argue by double induction: we shall assume that \(b_{i} \in \mathfrak{m}^{n-1}\) for all \(i\) and \(b_{i} \in \mathfrak{m}^{n}\) for \(i < k\) and show that this implies that \(b_{k} \in \mathfrak{m}^{n}\) . For this, we write \(f = \sum_{i=1}^{\infty} a_{i} X^{i}\) and compare the coefficients of \(\mathbf{X}^{s+k}\) in (3); then:

\[
\left(b _ {0} a _ {s + k} + \dots + b _ {k - 1} a _ {s + 1}\right) + b _ {k} a _ {s} + \left(b _ {k + 1} a _ {s - 1} + \dots + b _ {k + s} a _ {0}\right) = 0.\tag{4}
\]

The terms in the first bracket belong to \(\mathfrak{m}^n\) since the \(b_i \in \mathfrak{m}^n\) for \(i < k\) ; similarly for those in the second, since the \(b_i \in \mathfrak{m}^{n-1}\) for all \(i\) and the \(a_i \in \mathfrak{m}\) for \(i \leqslant s - 1\) . Hence \(b_k a_s \in \mathfrak{m}^n\) and, as \(a_s\) is an invertible element of A, \(b_k \in \mathfrak{m}^n\) .

\begin{enumerate}
\def\labelenumi{(\alph{enumi})}
\setcounter{enumi}{1}
\tightlist
\item
  We show that \(f\mathbf{B} + \mathbf{M} = \mathbf{B}\) . We write
\end{enumerate}

\[
g = a _ {s} + a _ {s + 1} \mathbf {X} + a _ {s + 2} \mathbf {X} ^ {2} + \dots ;
\]

it is an invertible element of B. Then

\[
f - \mathrm{X} ^ {s} g = a _ {0} + a _ {1} \mathrm{X} + \dots + a _ {s - 1} \mathrm{X} ^ {s - 1};
\]

if therefore we write \(fg^{-1} - \mathbf{X}^s = (f - \mathbf{X}^s g)g^{-1} = -h\) , the coefficients of \(h\) belong to \(m\) . Then let \(r\) be an element of B. By induction on \(n\) we define a sequence \((q^{(n)})\) of elements of B: we take \(q^{(0)}\) to be the unique series satisfying:

\[
r \equiv \mathbf {X} ^ {s} q ^ {(0)} \pmod {\mathbf {M}};\tag{5}
\]

writing \(h = \sum_{i=0}^{\infty} h_i X^i\) and \(q^{(n)} = \sum_{i=0}^{\infty} q_i^{(n)} X^i\) , the \(q_i^{(n)}\) are defined by:

\[
q _ {i} ^ {(n)} = \sum_ {j = 0} ^ {i + s} h _ {j} q _ {i + s - j} ^ {(n - 1)}.\tag{6}
\]

It follows immediately from (6) that:

\[
\mathbf {X} ^ {s} q ^ {(n)} \equiv h q ^ {(n - 1)} \pmod {\mathbf {M}}.\tag{7}
\]

As \(h_j \in \mathfrak{m}\) for all \(j\) , it also follows from (6), by induction on \(n\) , that \(q_i^{(n)} \in \mathfrak{m}^n\) for all \(i\) and all \(n\) . As \(A\) is complete, it follows that the series

\[
q ^ {(0)} + q ^ {(1)} + \dots + q ^ {(n)} + \dots
\]

converges to an element \(q\) of B. By (5) and (7),

\[
\mathbf {X} ^ {s} (q ^ {(0)} + q ^ {(1)} + \dots + q ^ {(n)}) \equiv r + h (q ^ {(0)} + \dots + q ^ {(n - 1)}) \pmod {\mathbf {M}}. \tag {8}
\]

As \(\mathbf{M}\) is closed, at the limit (8) gives \(r \equiv (\mathbf{X}^s - h)q\) (mod. \(\mathbf{M}\) ), that is

\[
r \in f g ^ {- 1} q + \mathbf {M} \subset f \mathbf {B} + \mathbf {M}.
\]

We may also use the results of Chapter III, §2 to show the relation \(B = fB + M\) (cf.~Exercise 12). The method followed here has the advantage of being applicable to convergent series.

COROLLARY. With the hypotheses and notation of Proposition 5, suppose that \(s \geqslant 1\) , so that \(f \in \mathbf{B}\mathfrak{m} + \mathbf{B}\mathbf{X}\) . Then the A-homomorphism \(h\) of \(\mathbf{B}' = \mathbf{A}[[\mathbf{T}]]\) to \(\mathbf{B} = \mathbf{A}[[\mathbf{X}]]\) such that \(h(\mathbf{T}) = f\) (Chapter III, §2, no. 9, Proposition 11 (a)) defines on \(\mathbf{B}\) a free \(\mathbf{B}'\) -module structure admitting \(\{1, \mathbf{X}, \ldots, \mathbf{X}^{s-1}\}\) as basis. In particular \(h\) is injective.

Let the B'-module B be given the (T)-adic filtration, which consists of the \(f^{n}B\) for \(n \geqslant 0\) (Chapter III, §2, no. 1). Then B/fB is a free module over the ring A = B'/TB' and the images of the \(X^{i}\) ( \(0 \leqslant i \leqslant s - 1\) ) in this A-module form a basis of it (Proposition 5); as moreover f is not a divisor of zero in B (Proposition 5), \(Bf^{n}/Bf^{n+1}\) is also a free (B'/TB')-module of rank s, so that condition (GR) of Chapter III, §2, no. 8 is satisfied (replacing A by B' and M by B). On the other hand, since B' is Hausdorff and complete with respect to the (T)-adic filtration and gr(B) is a finitely generated gr(B')-module by the above, it is seen first (Chapter III, §2, no. 9, Corollary 1 to Proposition 12) that B is a finitely generated B' module. The first assertion of the corollary then follows from Chapter III, §2, no. 9, Proposition 13. The second follows immediately from it.

DEFINITION 2. A polynomial \(\mathbf{F} \in \mathbf{A}[\mathbf{X}]\) is called distinguished if it is of the form

\[
\mathbf {F} = \mathbf {X} ^ {s} + a _ {s - 1} \mathbf {X} ^ {s - 1} + \dots + a _ {0},
\]

where \(a_{i} \in \mathfrak{m}\) for \(0 \leqslant i \leqslant s - 1\) .

Note that the product of two distinguished polynomials is a distinguished polynomial.

PROPOSITION 6 (Preparation Theorem). Let \(f \in \mathbf{B}\) be a series whose reduced series is not zero and \(s\) its reduced order. Then there exists a unique ordered pair \((u, \mathbf{F})\) such that \(u\) is an invertible element of \(\mathbf{B}\) , \(\mathbf{F}\) a distinguished polynomial of degree \(s\) and \(f = u\mathbf{F}\) .

We write \(F = X^{s} + G\) , where \(G = g_{0} + \cdots + g_{s-1}X^{s-1} (g_{i} \in A)\) . The relation f = uF is equivalent to \(F = u^{-1}f\) , that is to \(X^{s} = u^{-1}f - G\) . Hence Proposition 5 shows the uniqueness of G and \(u^{-1}\) and therefore of F and u. It also shows that there exist \(v \in B\) and a polynomial \(G = g_{0} + \cdots + g_{s-1}X^{s-1} (g_{i} \in A)\) such that \(X^{s} = vf - G\) ; it remains to show that v is invertible in B and that \(g_{i} \in m\) for all i. Now, writing \(\bar{g}_{i}\) for the canonical image of \(g_{i}\) in k and \(\bar{f}\) , \(\bar{v}\) for the reduced series of f, g,

\[
\mathbf {X} ^ {s} + \bar {g} _ {0} + \bar {g} _ {1} \mathbf {X} + \dots + \bar {g} _ {s - 1} \mathbf {X} ^ {s - 1} = \bar {f} \bar {v};
\]

since \(\bar{f}\) is of order \(s\) , \(\bar{g}_i = 0\) for all \(i\) and \(\bar{v}\) is of order 0, hence \(v\) is invertible.

PROPOSITION 7. Let F be a distinguished polynomial and g, h two formal power series of B such that F = gh. Then there exists an invertible element u of B such that ug and \(u^{-1}h\) are distinguished polynomials and \(\mathbf{F} = (ug)(u^{-1}h)\) .

In fact, the reduced series of g and h are \(\neq0\) ; hence, by Proposition 6, there exist invertible elements u, v of B such that ug and vh are distinguished polynomials. Then \(uvF = (ug)(vh)\) is a distinguished polynomial and uv is invertible. Passing to the reduced series, it is seen immediately that F and uvF have the same reduced order, that is the same degree. The uniqueness assertion in Proposition 6 therefore shows that F = uvF, whence uv = 1.

COROLLARY. Suppose further that A is an integral domain and F a distinguished polynomial of degree s. For F to be extremal in A{[}X{]}, it is necessary and sufficient that it be extremal in B = A{[}{[}X{]}{]}.

Suppose that F is not extremal in A{[}X{]}, so that \(F = f_{1}f_{2}\) , where \(f_{1}\) and \(f_{2}\) are non-invertible elements of A{[}X{]}; the product of the dominant coefficients of \(f_{1}\) and \(f_{2}\) being equal to 1, these coefficients are invertible in A and the hypothesis implies that \(f_{1}\) and \(f_{2}\) are of degrees \textgreater0 and \textless s; as the reduced polynomials \(\bar{f}_{1}, \bar{f}_{2}\) satisfy \(\bar{f}_{1}\bar{f}_{2} = X^{s}\) , neither \(\bar{f}_{1}\) nor \(\bar{f}_{2}\) can be invertible in k{[}{[}X{]}{]}, for, if \(\bar{f}_{1}\) were invertible, \(\bar{f}_{2}\) would be of order s, which is absurd. A fortiori, neither \(f_{1}\) nor \(f_{2}\) is invertible in B and F is not extremal in B.

Conversely, if F is not extremal in A{[}{[}X{]}{]}, then F = gh, where neither g nor h is invertible in B; their reduced orders are therefore \(\geqslant 1\) ; then the distinguished polynomials ug and \(u^{-1}h\) of Proposition 7 are not constant, which shows that F is not extremal in A{[}X{]}.

